%
% Phoenix-RTOS
%
% Documentation resources
%
% Look of the tables of the PDF documents
%
% Copyright 2026 Phoenix Systems
% Author: Damian Loewnau
%
% This file is part of Phoenix-RTOS.
%
% SPDX-License-Identifier: BSD-3-Clause

% The internals of sphinx and of the array package are used below
\makeatletter

% Rules: a thin grey grid with the orange rule below the header row
\definecolor{ps-tableline}{HTML}{BFBFBF}

\newlength{\psTableRuleWidth}
\newlength{\psTableHeaderRuleWidth}
\setlength{\psTableRuleWidth}{0.8pt}
\setlength{\psTableHeaderRuleWidth}{2pt}

\arrayrulecolor{ps-tableline}
\setlength{\arrayrulewidth}{\psTableRuleWidth}

% Least width of a column of a table sized by tabulary, which otherwise shrinks
% every column in the same proportion and can leave a short header clipped
\setlength{\tymin}{3em}

% Space kept above and below the content of a cell, the header row getting more
% of it. The space is part of the cell, so that the vertical rules span it.
\newlength{\psTableCellPadding}
\newlength{\psTableHeaderPadding}
\setlength{\psTableCellPadding}{11pt}
\setlength{\psTableHeaderPadding}{6pt}

% A table standing in a box, an admonition among them, keeps a tighter padding:
% such a box does not break across pages, so the taller rows would overflow it.
% The flag is taken before sphinx marks the table itself as a box, which it does
% for the sake of the code blocks that may sit in its cells.
\newif\ifpsTableBoxed
\newlength{\psTableBoxedCellPadding}
\setlength{\psTableBoxedCellPadding}{3pt}

\newcommand{\psTableCellPaddingNow}{%
	\ifpsTableBoxed\psTableBoxedCellPadding\else\psTableCellPadding\fi}

\let\psTableStart\sphinxattablestart

\def\sphinxattablestart{%
	\ifspx@inframed\psTableBoxedtrue\else\psTableBoxedfalse\fi
	\psTableStart
}

% Space between the table and its caption
\newlength{\psTableCaptionSkip}
\setlength{\psTableCaptionSkip}{4pt}

% Struts keeping the given space above and below the line they are put in
\newcommand{\psTableStrutAbove}[1]{%
	\rule{\z@}{\dimexpr\ht\strutbox+#1\relax}}
\newcommand{\psTableStrutBelow}[1]{%
	\rule[\dimexpr-\dp\strutbox-#1\relax]{\z@}{\dimexpr\dp\strutbox+#1\relax}}

\newcommand{\psTableHeaderCellPadding}{%
	\dimexpr\psTableCellPaddingNow+\psTableHeaderPadding\relax}

% Padding of a cell, put around its content by the column types below. A header
% cell asks for more of it by the flag, which lives in the group of that single
% cell, hence it needs no cleaning up.
\newif\ifpsTableHeaderCell

\newcommand{\psTableCellTop}{\psTableStrutAbove{\psTableCellPaddingNow}}
\newcommand{\psTableCellBottom}{%
	\ifpsTableHeaderCell
		\psTableStrutBelow{\psTableHeaderCellPadding}%
	\else
		\psTableStrutBelow{\psTableCellPaddingNow}%
	\fi}

% Width of a column given as the #1/#2 fraction of the table width, copied from
% the \X column type of sphinx (see sphinxlatextables.sty)
\newcommand{\psTableColumnWidth}[2]{\dimexpr
	(\linewidth-\spx@arrayrulewidth)*#1/#2-\tw@\tabcolsep-\spx@arrayrulewidth\relax}

% The column types of sphinx, redefined to pad their cells. T is used for the
% tables of an automatic layout, \X and \Y for the ones of given widths, the
% width of \Y being copied from sphinx as well.
\newcolumntype{T}{>{\psTableCellTop}J<{\psTableCellBottom}}
\newcolumntype{\X}[2]{%
	>{\psTableCellTop}p{\psTableColumnWidth{#1}{#2}}<{\psTableCellBottom}}
\newcolumntype{\Y}[1]{%
	>{\psTableCellTop}p{\dimexpr
		#1\dimexpr\linewidth-\spx@arrayrulewidth\relax
		-\tw@\tabcolsep-\spx@arrayrulewidth\relax}<{\psTableCellBottom}}

% Cells of a data table, centered both horizontally and vertically, as the
% column types of sphinx are all aligned to the top
\newcolumntype{\psC}[2]{%
	>{\centering\arraybackslash\psTableCellTop}m{\psTableColumnWidth{#1}{#2}}%
	<{\psTableCellBottom}}

% Header cells are bold, centered and taller than the ones of the body. The
% strut above has to repeat the padding of the cell, as only the tallest strut
% of a line counts, while the space below is left to \psTableCellBottom, which
% is put after the last line of the cell.
\protected\def\sphinxstyletheadfamily{%
	\sffamily\bfseries\centering\arraybackslash
	\psTableHeaderCelltrue
	\psTableStrutAbove{\psTableHeaderCellPadding}}

% The rules are drawn on their own instead of by \hline, which longtable splits
% into two rules with a page break in between, one of them ending a page and the
% other one starting the next one. Such a leftover rule shows up right below the
% orange rule of the head repeated on that page.
\newcommand{\psTableRule}[2]{\noalign{{\color{#1}\hrule height#2}}}

% Drawing the rules takes those page breaks away, so a row rule brings one back,
% after the rule, which keeps the rule on the page of the row it closes
\newcommand{\psTableRowRule}{%
	\psTableRule{ps-tableline}{\psTableRuleWidth}%
	\noalign{\penalty-\@lowpenalty}%
}

\newcommand{\psTableHeaderRule}{\psTableRule{ps-orange}{\psTableHeaderRuleWidth}}

% Sphinx puts the caption above the table while the documents want it below, so
% the caption is only remembered at that point and typeset once the table ends.
% Of the caption, only its label, \fnum@table, follows the language of the
% document, while the text is taken as it was written.
\let\psSavedTableCaption\relax

\newcommand{\psRememberTableCaption}[2]{%
	\addcontentsline{lot}{table}{\protect\numberline{\thetable}{\ignorespaces #1}}%
	\gdef\psSavedTableCaption{#2}%
}

% A table of a single page has nothing else to step its counter, and it has to be
% stepped before the \label that sphinx puts right after the caption
\long\def\psSaveTableCaption[#1]#2{%
	\refstepcounter{table}%
	\psRememberTableCaption{#1}{#2}%
}

\newcommand{\psTypesetTableCaption}{%
	\ifx\psSavedTableCaption\relax\else
		\par\nobreak
		\vskip\psTableCaptionSkip
		\noindent\vtop{\hsize\linewidth\raggedleft
			\fnum@table: \psSavedTableCaption\par}%
		\par
		\global\let\psSavedTableCaption\relax
	\fi
}

% The caption of a multi-page table is turned by longtable into a row of the
% head, which would keep it above the table. Sphinx emits it as
%     \caption{...}\label{...}\\*[\sphinxlongtablecapskipadjust]
% so the row that terminates it and the space adjusting it are swallowed here as
% well, leaving the caption itself to \psTypesetTableCaption below the table.
% The pattern is the one of the table templates of the pinned sphinx.
\long\def\psTakeOverLongtableCaption#1\\*[#2]{\noalign{\psGrabLongtableCaption#1}}

% \LT@array of longtable has stepped the counter at the start of the table
\long\def\psGrabLongtableCaption#1{\psRememberTableCaption{#1}{#1}}

\let\psLongtableStart\sphinxatlongtablestart

\def\sphinxatlongtablestart{%
	\ifspx@inframed\psTableBoxedtrue\else\psTableBoxedfalse\fi
	\psLongtableStart
	% \caption of a longtable is \LT@caption, which longtable itself installs
	\let\LT@caption\psTakeOverLongtableCaption
}

% All of the above is applied to every table by extending the style macro that
% sphinx executes inside the scope of a single table
\let\psTableStandardStyle\sphinxthistablewithstandardstyle

\def\sphinxthistablewithstandardstyle{%
	\psTableStandardStyle
	\let\sphinxmidrule\psTableHeaderRule
	\let\sphinxhline\psTableRowRule
	\let\spx@toprule\psTableRowRule
	\let\sphinxbottomrule\psTableRowRule
	\def\sphinxcaption{\@dblarg\psSaveTableCaption}%
	% the caption is gone from above the table, so undo the space made for it
	\def\sphinxaftertopcaption{\vskip\baselineskip\vskip\parskip}%
	\let\sphinxtableafterendhook\psTypesetTableCaption
}

\makeatother
